% !TeX root = ../papers/nc-p3.tex
\section{可変跡の円板での六次中心商の二文字の確定}
\label{sec:rank-fifteen-saturated-center-products}
\begin{proposition}
定理\ref{thm:rank-fifteen-variable-zero-disk}の領域で前節の四中心式による商を考える。
その長さ6の文字別次元を $d_h$ とすると、全始点で
\[
            d_0=d_1=12,\qquad 10\le d_2,d_3\le12 .
\]
従って六次中心商の総次元は44以上48以下である。
\end{proposition}
\begin{proof}
$I$ を四中心式の生成するイデアルとする。
前節の各四行十四列の行列を精度9から計算し、単元による行操作と
$H$ による飽和で得た階数4の行空間を、主列が単位行列になる基底で表す。
残り一桁の精度でその基底の全定数項が得られる。飽和の前後で $H\ne0$ に
おける行空間は同じなので、この基底は長さ5の $I$ の元からなる。
これを四つの末尾文字に沿って直和し、長さ6の未切断経路空間へ射影する。
この射影は前節までと同じ第六挿入行列で計算する。
五桁のうち二桁の損失で求まるので、積の基底・射影とも定数項が使える。
四式は中心的で経路は一次生成なので $I_6=I_5 A_1$ であり、
各文字について得る十六行は $I$ の六次成分を張る。

\path{computations/verify_rank_fifteen_saturated_center_products.py}
は、これらの定数項を $z=1+i$、標数3に還元して消去する。
四始点すべてで、十六行の階数は文字順に $(8,8,10,10)$ である。
対応する小行列式は単元なので、これは円板全体での $I$ の次数6成分の次元の下界を与える。
一方、中心性により各始点の六次成分は、中心式と長さ2の経路の積で生成される。
長さ2の文字別次元が $(2,2,3,3)$、中心式は文字0と文字1が各二本なので、
生成元数は文字ごとに $(8,8,12,12)$ である。これらを未切断経路の次元
$(20,20,22,22)$ から引けば主張を得る。

さらに文字 $h=2,3$ では、十個の単元主元を用いて十六行二十二列の行列を
十次の単位行列と六行十二列の残りの行列 $B_h$ に分けられる。
したがって $d_h=12-\operatorname{rank} B_h$、かつ $\operatorname{rank} B_h\le2$ である。
この文字で目標次元10に達する必要十分条件は $B_h$ の二次小行列式の少なくとも一つが
その点で非零となることである。
\end{proof}
